\subsection{A metric binary-tree obstruction and a fourth quotient class}

\begin{claim}{A.30}{Proved}{A countable diagonal block from an invariant length}
Let $E$ be a countable nonabelian simple group with a function
$\ell:E\to[0,\infty)$ satisfying
\[
 \begin{gathered}
 \ell(1)=0,\quad \ell(g)>0\ (g\ne1),\quad
 \ell(g^{-1})=\ell(g),
 \\
 \ell(gh)\leq\ell(g)+\ell(h),\qquad
 \ell(aga^{-1})=\ell(g),\qquad
 \inf_{g\ne1}\ell(g)=0.
 \end{gathered}
 \tag{A.30.1}
\]
Let $B$ be any specified connected CW classifying space with
fundamental group $H$, let $K\leq H$, and let
$q:K\twoheadrightarrow E$ be an epimorphism.
There is an actual inclusion $B\hookrightarrow Y$ obtained by
attaching countably many one-cells and two-cells such that:
\begin{enumerate}
\item $Y$ is aspherical, $\dim Y\leq\max\{\dim B,2\}$, and
$H\to P:=\pi_1Y$ is injective;
\item $B\to Y$ is an ordinary integral homology isomorphism,
and the displayed copy of $H$ normally generates $P$;
\item no homomorphism $P\to E$ restricts to $q$ on the
specified subgroup $K$.
\end{enumerate}
The relative cellular boundary from the attached two-cells to
the attached one-cells is the identity in the specified bases.
No map from all of $H$ to $E$ is required.

For every starting space $B$ in A.28, a new construction
$B\hookrightarrow X^{\sharp}$ has all the conclusions of A.28
and also excludes countable nonabelian simple quotients admitting
a length as in (A.30.1). In particular there is an actual
two-dimensional acyclic map
\[
 t^{\sharp}:BG^{\sharp}\longrightarrow S^2,
 \qquad |G^{\sharp}|\leq 2^{\aleph_0},
 \tag{A.30.2}
\]
with the three quotient exclusions, equation-solving property,
and torsion and locally linear image exclusions of A.28, and
with this fourth quotient exclusion. Its specified A.10
resolution is an actual subcomplex, and its map extends the
original resolution map literally.
\end{claim}

\begin{proof}
Put $d(g,h)=\ell(g^{-1}h)$. The displayed axioms make this a
bi-invariant metric. We use the convention
$[a,b]=aba^{-1}b^{-1}$.

\paragraph{A small symmetric normal generating set.}
For $b\ne1$, define
\[
 T_b(u,v)=v[b,u]v^{-1},\qquad
 S_b=\{T_b(u,v):u,v\in E\}.
 \tag{A.30.3}
\]
The set $S_b$ contains $1$ and is invariant under conjugation.
It is symmetric, because
\[
 [b,u]^{-1}=u[b,u^{-1}]u^{-1},\qquad
 T_b(u,v)^{-1}=T_b(u^{-1},vu).
 \tag{A.30.4}
\]
The center of a nonabelian simple group is trivial, so some
$[b,u]$ is nonidentity. Consequently the subgroup generated
by $S_b$ is a nontrivial normal subgroup and equals $E$.
The length bound needed below is
\[
 \ell(s)\leq2\ell(b)\qquad(s\in S_b).
 \tag{A.30.5}
\]
Indeed $[b,u]=b(ub^{-1}u^{-1})$, and conjugation preserves
length. No multiplicative contraction estimate for
commutators is needed.

\paragraph{A finite word that stays nonconstant while losing one value.}
Let $R:E^m\to E$ be a nonconstant word map with finitely many
coefficients. Give $E^m$ the maximum of its coordinate metrics.
If a displayed word for $R$ has $L$ occurrences of variables
or their inverses, the triangle inequality and bi-invariance
give
\[
 d(R(\mathbf x),R(\mathbf y))
 \leq L\max_i d(x_i,y_i).
 \tag{A.30.6}
\]
This follows by replacing its variable occurrences one at a
time; fixed left and right multiplication, and inversion,
are isometries. Nonconstancy ensures that $L\geq1$.

Fix a value $e\in E$. Choose $\mathbf a\in E^m$ with
$c:=R(\mathbf a)\ne e$, and write $\delta=d(c,e)>0$.
Choose $b\ne1$ so small that
\[
 2L\ell(b)<\delta.
 \tag{A.30.7}
\]
This uses exactly the last condition in (A.30.1).
Choose also $\mathbf a'$ with $R(\mathbf a')\ne c$.
Since $S_b$ generates $E$ and contains $1$, its Cartesian
power $S_b^m$ generates $E^m$. More explicitly, express each
coordinate of $\mathbf a'\mathbf a^{-1}$ as a finite product
of elements of $S_b$, and pad these products by $1$ to have
equal lengths. Thus there is a finite sequence from
$\mathbf a$ to $\mathbf a'$ whose successive steps are left
multiplications by tuples in $S_b^m$.

Take its first step at which the value of $R$ changes from
$c$. Its endpoints have the form
\[
 \mathbf u,\quad\mathbf s\mathbf u,
 \qquad \mathbf s\in S_b^m,
 \qquad R(\mathbf u)=c\ne R(\mathbf s\mathbf u).
 \tag{A.30.8}
\]
For every $\mathbf t\in S_b^m$, (A.30.5)--(A.30.7) imply
\[
 d(R(\mathbf t\mathbf u),c)
 \leq L\max_i\ell(t_i)
 \leq 2L\ell(b)<\delta.
 \tag{A.30.9}
\]
Here $d(t_i u_i,u_i)=\ell(t_i)$ uses conjugation invariance.
It follows that $R(S_b^m\mathbf u)$ excludes $e$ and contains
at least two elements, the two values in (A.30.8).

Replace each coordinate variable $x_i$ by
$T_b(y_i,z_i)u_i$, using a separate pair $(y_i,z_i)$ for
each $i$. The image of the resulting word map is exactly
$R(S_b^m\mathbf u)$: the definition of $S_b$ gives both
inclusions, and the separate pairs make the coordinate
choices independent. We have therefore produced a new
nonconstant finite word map, with image contained in the
old image and avoiding the prescribed value $e$.

\paragraph{The actual countable system.}
Enumerate $E$ as $e_0,e_1,\ldots$. The group is infinite,
since a positive definite length on a finite group cannot
have the last property in (A.30.1). Let
\[
 \mathcal T=\bigcup_{n\geq0}\{0,1\}^n
\]
be the finite binary words, including the empty word.
Start with $R_0(x_{\varnothing})=x_{\varnothing}$.
At step $n$ apply the preceding construction to the
nonconstant word $R_n$ on the $2^n$ variables indexed by
words $s$ of length $n$, and to $e_n$. It gives a coefficient
$b_n\ne1$ and a tuple $(u_s)_{|s|=n}$.
Define $R_{n+1}$ by substituting, simultaneously,
\[
 x_s=T_{b_n}(x_{s0},x_{s1})u_s
 \qquad(|s|=n).
 \tag{A.30.10}
\]
Writing $I_n=R_n(E^{2^n})$, the construction proves
\[
 I_{n+1}\subseteq I_n,\qquad |I_{n+1}|\geq2,
 \qquad e_n\notin I_{n+1}.
 \tag{A.30.11}
\]
Take (A.30.10) for every $s\in\mathcal T$ as one
countable system. A solution in $E$ would have
$x_{\varnothing}=R_n((x_s)_{|s|=n})\in I_n$ for every
$n$, by finite substitution. But each element $e_n$ is
excluded from $I_{n+1}$. There is no solution.

\paragraph{Lifting coefficients and retaining actual classifying spaces.}
Choose lifts $\beta_n\in K$ of $b_n$ and
$\upsilon_s\in K$ of $u_s$, and choose based loops in
$B$ representing these lifts. The same selected lift is
used at every occurrence of each selected coefficient.
Attach a loop $x_s$ for each $s\in\mathcal T$ and a
two-cell with attaching word
\[
 r_s=x_s\upsilon_s^{-1}
 \bigl(x_{s1}[\beta_{|s|},x_{s0}]x_{s1}^{-1}\bigr)^{-1}.
 \tag{A.30.12}
\]
Let $Y$ be the resulting actual CW complex. Every child
variable has total exponent zero in $r_s$, and the parent
$x_s$ has exponent one. The relative cellular chain map is
therefore exactly
\[
 C_2(Y,B)=\bigoplus_{s\in\mathcal T}\mathbb Z[r_s]
 \longrightarrow
 C_1(Y,B)=\bigoplus_{s\in\mathcal T}\mathbb Z[x_s],
 \qquad [r_s]\longmapsto[x_s].
 \tag{A.30.13}
\]
There are no other relative cells. This proves the ordinary
integral homology assertion.

For the asphericity verification, let $Y_n$ consist of $B$,
the loops with $|s|\leq n$, and the two-cells with $|s|<n$.
In $Y_n$, the root loop occurs once in its own attaching
word and occurs in no other attaching word. Its loop and
two-cell can be removed by the elementary presentation
collapse across that two-cell, fixing the rest of the
complex. The replacement path for the removed loop is
precisely the right side of (A.30.10), with lifted
coefficients. After this removal, the same argument applies
to each loop at depth one, then at depth two, and so on.
In increasing depth order these are actual homotopy
equivalences relative to $B$ and the frontier loops, giving
\[
 Y_n\simeq_B B\vee\bigvee_{|s|=n}S^1.
 \tag{A.30.14}
\]
At each removal the parent appears in no remaining
two-cell other than its own: its parent's two-cell has
already been removed, and all other coefficients lie in
$B$. This verifies the required free occurrence, without
an infinite presentation elimination.

The space on the right of (A.30.14) is aspherical. Its
universal cover is a tree of contractible lifted copies
of $B$ connected by intervals. Also $H\to\pi_1Y_n$ is
the inclusion of the displayed free factor. The spaces
$Y_n$ form actual CW subcomplexes with union $Y$.
A map from a sphere of dimension at least two has image
in a finite subcomplex and hence in some $Y_n$; it is
nullhomotopic there. Likewise a nullhomotopy in $Y$ of a
loop from $B$ is supported in some $Y_n$, where the
free-factor inclusion is injective. Thus $Y$ is
aspherical and $H\to\pi_1Y$ is injective. No injectivity
of all transition homomorphisms $\pi_1Y_n\to\pi_1Y_{n+1}$
was used.

If the image of $H$ is killed in $P=\pi_1Y$, every
$\beta_n$ and $\upsilon_s$ becomes $1$. Each relation
(A.30.12) then reads $x_s=1$. This proves normal
generation by $H$. Finally a homomorphism $P\to E$
restricting to $q$ on $K$ would send all these relations
to (A.30.10), producing the excluded solution. The block
assertions, including their specified comparison maps,
are proved.

\paragraph{Rerunning the construction with four classes.}
Let $\mathcal C^{\sharp}$ be the union of the three
classes in A.28 and the class of countable nonabelian
simple groups equipped with a length as in (A.30.1).
We describe a new construction, denoting its final
group by $G^{\sharp}$.

Start with the specified countable $B$ of A.28 and run
$\omega_1$ passes. At the beginning of a pass, with
group $A$ of cardinality at most $\mathfrak c$, list
every triple
\[
 K\leq A\text{ countable},\qquad
 E\in\mathcal C^{\sharp},\qquad
 q:K\twoheadrightarrow E.
 \tag{A.30.15}
\]
Use representatives on countable underlying sets.
There are at most $\mathfrak c$ such triples: there
are at most $\mathfrak c^{\aleph_0}=\mathfrak c$
countable subgroups and maps, and at most
$\mathfrak c$ group structures on a countable set.
The functions $\ell$ on a countable group also number
at most $\mathfrak c^{\aleph_0}=\mathfrak c$.
One can either include these witnesses in the list or
choose one for each eligible group.

Process the triples in a well-order of cardinality at
most $\mathfrak c$. For the first three classes attach
the subgroup block of A.27, as in A.28. For the new
class attach (A.30.12) over the entire current
classifying space, lifting coefficients from the
specified $K$. The preceding proof applies without
requiring any homomorphism from that entire group
to $E$. Earlier subgroups persist by the proved
coefficient injections.

In each pass also process all one-variable words of
total exponent $1$ or $-1$ over its starting group,
using the one-generator, one-relator adjunction in
A.18. There are at most $\mathfrak c$ such words.
As in A.28, all intermediate two-dimensional
classifying groups are torsion-free, so that the
audited Forester--Rourke adjunction applies. Each
such pair has relative cellular boundary $\pm1$
and normally generates no new quotient after its
coefficient group is killed.

Every block has countably many cells. Each pass,
and the union of $\omega_1$ passes, has at most
$\mathfrak c$ cells. At limit stages use the
actual CW union. The finite-support argument for
spheres and nullhomotopies proves asphericity and
preserves the initial group injection; filtered
cellular chain unions preserve the integral
homology isomorphisms. Normal generation is
transitive under the successive attachments and
persists at unions. Consequently the new inclusion
$B\hookrightarrow X^{\sharp}$ has exactly the
dimension, cardinality, group injection, homology,
and normal-generation properties stated in A.28.

Every one-variable unimodular equation over
$G^{\sharp}$ uses finitely many coefficients,
all present in one earlier pass. Its word is
therefore processed subsequently. Any quotient
inherits solutions by lifting its finite list of
coefficients and mapping a solution down. Thus
the equation-solving conclusion of A.28 is
retained as well.

Suppose $G^{\sharp}\twoheadrightarrow E$ were a
nontrivial countable quotient in one of the four
classes. Choose one preimage for each element of
$E$. All these countably many preimages belong
to one earlier outer stage, since $\omega_1$ is
regular. They generate a countable subgroup $K$
whose restriction maps onto $E$. The corresponding
triple (A.30.15), with a chosen invariant length
in the fourth case, is processed in a subsequent
pass. Its block forbids the restriction of the
assumed final homomorphism. This contradiction
proves all four exclusions.

For the specified finite A.10 model over $S^2$,
normal generation transfers its absence of finite,
torsion, and locally linear images exactly through
the homomorphisms proved in A.28: a map from
$G^{\sharp}$ to any such target kills every
inserted Higman group, then kills $H$, then kills
$G^{\sharp}$. Extend the original map over each
new loop by the constant loop, and over each
two-cell by a nullhomotopy of its image in the
simply connected sphere. CW unions give an actual
map with $t^{\sharp}|_B=t_0$. Its integral homology
map is an isomorphism because the inclusion and
$t_0$ have that property. All local systems on
$S^2$ are constant, so the universal coefficient
theorem gives the acyclicity required in A.2.
Its nonzero second homology makes the dimension
exactly two. This proves (A.30.2).
\end{proof}

For this new $S^2$ construction, finite normal
generation gives the same infinite simple quotient
reduction as in A.28. A remaining countable simple
quotient solves all one-variable unimodular
equations, so A.29 supplies a nondiscrete
Hausdorff metrizable group topology. At the same
time, A.30 excludes every length on that quotient
satisfying all the conditions (A.30.1). The
conjugation estimate proved for the metric in
A.29 is $L(ghg^{-1})\leq D_gL(h)$; the argument
above requires equality under conjugation in
order to obtain the uniform bound (A.30.5).
